\documentclass[10pt,a4paper]{amsart}
\usepackage[margin=19mm]{geometry}
\usepackage{amsmath,amssymb,mathtools}
\usepackage[scr=rsfs,cal=euler]{mathalfa}
\usepackage{xfrac}
\usepackage[unicode,hidelinks,bookmarks=false]{hyperref}
\newcommand{\ie}{i.e.\ }
\newcommand{\cf}{cf.\ }
\newcommand{\resp}{respectively\ }
\newcommand{\Subsection}{\subsection}
\providecommand{\nobreakdash}{\nobreak\hbox{-}}
\setlength{\emergencystretch}{2em}
\multlinegap0pt
\usepackage{amssymb}
\usepackage[shortlabels]{enumitem}
\newtheorem{proposition}{Proposition}
\newcommand{\bigmid}{\,\middle\vert\,}
\renewcommand{\d}{\mathrm{d}}
\title{Rotationally symmetric critical metrics for Laplace eigenvalues on tori in a conformal class}
\author{Egor Morozov}
\date{Source: arXiv:2601.19196v1 (CC BY 4.0)}
\begin{document}
\maketitle

\noindent\emph{Source and licence.} Rotationally symmetric critical metrics for Laplace eigenvalues on tori in a conformal class. Version arXiv:2601.19196v1 (CC BY 4.0), distributed by arXiv under the Creative Commons Attribution 4.0 International licence, http://creativecommons.org/licenses/by/4.0/, as recorded in the arXiv licence metadata for this exact version and shown on the abstract page https://arxiv.org/abs/2601.19196v1. Full licence notice: \texttt{licenses/arxiv-2601.19196-CC-BY-4.0.txt}.\par\smallskip

\noindent\textit{Editorial scope.} Counted unit: the article's proposition pr:weyl-N2eqscsc (source0.tex lines 1247--1250). The article's complete original proof of it is delivered with it (source0.tex lines 1252--1355, one \texttt{proof} environment of 104 lines). No mathematics is written, repaired or re-derived: the statement and the proof are the article's own lines, in the article's own notation, and no retained line is substituted or re-flowed.  Retained uncounted as the source-local context the proof needs: lines 417--419; lines 421--425; lines 599--602; lines 628--632.  Nothing was omitted from any retained span, so the package carries no omission record.  No retained line contains a citation, so the wrapper carries no bibliography and no bibliography entry.  No retained line contains a figure environment, an image-inclusion command or drawing code, so no image is delivered and the package declares no asset dependency.  Editorial formatting only, in addition: an AMS \texttt{amsart} preamble that restates the article's own declarations and macros used by the retained lines, namely the environments \texttt{proposition} on separate counters, together with the standard packages those lines need.  The retained environments print in this document as Proposition 1 in order of appearance, whereas the article prints the same items with the numbering of its own full document.  The complete native source of the article as distributed in the arXiv e-print for this exact version is bundled unchanged as \texttt{sources/193441/source0.tex}.
\begin{equation}\label{eq:intro-N2eqsc}
\tau_2+\tau_3-\tau_1\le 4.
\end{equation}

\begin{equation}\label{eq:intro-N2eqscsc}
\text{either}\quad
\frac{p}q>\frac{1}{\sqrt 3}\quad\text{or}\quad
\left|\frac{r+a}q\right|<\frac{\sqrt 3}4.
\end{equation}

\begin{equation}\label{eq:taus-n0n1}
\Phi(n_0\mid m)=\frac{\pi p}q,\qquad
\Phi(n_1\mid m)=\left|\frac{\pi (r+a)}q\right|.
\end{equation}

\begin{equation}\label{eq:taus-limzero}
\lim_{m\to 0}\frac{m}{n_0}=1-\frac{4p^2}{q^2},\qquad
\lim_{m\to 0}\frac{m}{n_1}=1-\frac{4(r+a)^2}{q^2},\qquad
\lim_{m\to 0}\Psi(m)=\frac{\pi^2}{q^2}(p^2-(r+a)^2),
\end{equation}

\begin{proposition}\label{pr:weyl-N2eqscsc}
If~\eqref{eq:intro-N2eqscsc} is satisfied, then~\eqref{eq:intro-N2eqsc} is
satisfied.
\end{proposition}

\begin{proof}
Rewriting the LHS of~\eqref{eq:intro-N2eqsc} in terms of \(n_0,n_1,m\) we obtain
\[
\tau_2+\tau_3-\tau_1=\frac{n_1}{n_1-n_0}\left(1-n_0-\frac{n_0}m\right).
\]
Consider two cases.

\textit{Case 1: \(\frac{p}q>\frac{1}{\sqrt 3}\).} We have
\[
\frac{n_1}{n_1-n_0}\left(1-n_0-\frac{n_0}m\right)\le
\frac{1}{1-n_0}\left(1-n_0-\frac{n_0}m\right)=
1-\frac{n_0}{(1-n_0)m}.
\]
Hence it suffices to show that \(\Theta(m)>0\), where
\[
\Theta(m):=3+\frac{n_0}{(1-n_0)m}
\]
and \(n_0\) is a function in \(m\) given implicitly by~\eqref{eq:taus-n0n1}. It
follows from~\eqref{eq:taus-limzero} that
\[
\lim_{m\to 0}\Theta(m)=3+\frac{1}{1-\frac{4p^2}{q^2}}>0,\quad\text{since}\quad
\frac{p}q>\frac{1}{\sqrt 3}.
\]
We have to show that \(\Theta(m)\) does not change sign on \((0,1)\). For this
it suffices to show that if \(\Theta(m)=0\) for some \(m\in (0,1)\), then
\[
\Phi(n_0\mid m)=\sqrt{\frac{(1-n_0)(n_0-m)}{n_0}}\Pi(n_0\mid m)<\frac{p\pi}q.
\]
Eliminating \(n_0\) and using again that \(\frac{p}q>\frac{1}{\sqrt 3}\), we
obtain that it suffices to prove the following inequality
\[
\Omega(m):=\sqrt{\frac{4-3m}{3(1-3m)}}\Pi\left(-\frac{3m}{1-3m}\bigmid
m\right)\le\frac{\pi}{\sqrt 3},\qquad m\in (0,1/3).
\]
We have
\[
\Omega'(m)=\frac{2(2-3m)E(m)-(4-3m)(1-m)K(m)}{2m(1-m)\sqrt{3(4-3m)(1-3m)}}.
\]
Since \(\Omega(0)=\frac{\pi}{\sqrt 3}\), it remains to check that
\[
2(2-3m)E(m)-(4-3m)(1-m)K(m)<0,\qquad m\in (0,1/3).
\]
We have
\[
\frac{K(m)}{E(m)}>1+\frac{m}4>\frac{2(2-3m)}{(4-3m)(1-m)},
\]
where the first inequality follows from \(K(0)=E(0)=\pi/2\) and
\[
\frac{\d}{\d m}\left[K(m)-\left(1+\frac{m}4\right)E(m)\right]=
\frac{1}8\left(\frac{(1+3m)E(m)}{1-m}+K(m)\right)>0,
\]
and the second inequality is equivalent to the trivial inequality
\(m^2(3m+5)>0\).

\textit{Case 2: \(\bigl|\frac{r+a}q\bigr|<\frac{\sqrt 3}4\).} It is easy to see
that
\[
\frac{n_1}{n_1-n_0}\left(1-n_0-\frac{n_0}m\right)\le
n_1\left(1+\frac{1}m\right).
\]
Hence it suffices to show that \(\Theta(m)>0\), where
\[
\Theta(m):=4-n_1\left(1+\frac{1}m\right).
\]
and \(n_1\) is a function in \(m\) given implicitly by~\eqref{eq:taus-n0n1}. It
follows from~\eqref{eq:taus-limzero} that
\[
\lim_{m\to 0}\Theta(m)=4-\frac{1}{1-\frac{4(r+a)^2}{q^2}}>0,\quad\text{since}\quad
\left|\frac{r+a}q\right|<\frac{\sqrt 3}4.
\]
We have to show that \(\Theta(m)\) does not change sign on \((0,1)\). For this
it suffices to show that if \(\Theta(m)=0\) for some \(m\in (0,1)\), then
\[
\Phi(n_1\mid m)=\sqrt{\frac{(1-n_1)(n_1-m)}{n_1}}\Pi(n_1\mid m)>
\left|\frac{(r+a)\pi}q\right|.
\]
Eliminating \(n_1\) and using again that \(\bigl|\frac{r+a}q\bigr|<\frac{\sqrt
3}4\), we obtain that it suffices to prove the following inequality
\[
\Omega(m):=\sqrt{\frac{(1-3m)(3-m)}{4(1+m)}}\Pi\left(\frac{4m}{1+m}\bigmid
m\right)\ge\frac{\sqrt 3}4\pi,\qquad m\in (0,1/3).
\]
We have
\[
\Omega'(m)=\frac{(3m^2-2m+3)E(m)-(3-m)(1-m)K(m)}{4m(1-m)\sqrt{(3-m)(1-3m)(1+m)}}.
\]
Since \(\Omega(0)=\frac{\sqrt 3\pi}4\), it remains to check that
\[
(3m^2-2m+3)E(m)-(3-m)(1-m)K(m)>0,\qquad m\in (0,1/3).
\]
We have
\[
\frac{K(m)}{E(m)}<\frac{1}{\sqrt{1-m}}<\frac{3m^2-2m+3}{(3-m)(1-m)},
\]
where the first inequality follows from \(K(0)=E(0)=\pi/2\) and
\[
\frac{\d}{\d m}\left[\sqrt{1-m}K(m)-E(m)\right]=
\frac{(1-\sqrt{1-m})(E(m)-K(m))}{2m\sqrt{1-m}}<0,
\]
and the second inequality is equivalent to the trivial inequality
\[
m(9m^3-11m^2+15m+3)>0,\qquad m\in (0,1/3).
\]
\end{proof}

\end{document}
